\section{Introduction}
\label{sec:introduction}

The Track R program (R.0--R.3) has established that all gaming vectors
available to a hostile actor deploying the \bisect{} redistricting platform
are either bounded below the noise floor (R.1: parameter gaming, $\leq 0.3$
seats; R.2: input data manipulation, $\leq 0.4$ seats) or pre-specified
by statute (R.3: geographic definitions).

But bounding the maximum effect of a gaming attack is not sufficient
for legal purposes.
A court needs to know: given this specific plan, can we verify that no
gaming occurred?
This is a question of \emph{provenance}---the ability to trace the plan
from its inputs to its outputs and verify that the computation was performed
as claimed.

Provenance has a long history in forensic science and litigation.
The chain of custody for physical evidence establishes provenance: who
had the evidence, when, and what they did with it.
For algorithmic redistricting, the analogous chain of custody is
cryptographic: SHA-256 hashes link the Census Bureau's public data releases
through the computation to the final plan, with each link in the chain
independently verifiable.

This paper formalises the \bisect{} audit chain as a legal defense.
We address three questions:
\begin{enumerate}
\item \textbf{Cryptographic}: How strong is the SHA-256 guarantee?
  What is the probability that an adversary can tamper with any link
  in the chain without detection?
  (Section~\ref{sec:sha256})
\item \textbf{Technical}: How does the three-level provenance architecture
  work in practice, and how does \texttt{bisect label-verify} demonstrate
  it for a specific plan?
  (Sections~\ref{sec:provenance}--\ref{sec:demo})
\item \textbf{Legal}: Does the audit chain satisfy the Daubert standard
  for admissibility of scientific evidence?
  What statutory language should a state legislature adopt to make the
  manifest filing enforceable?
  (Sections~\ref{sec:daubert}--\ref{sec:legal})
\end{enumerate}

\subsection{Why Reproducibility Matters in Redistricting Litigation}

Expert witnesses in redistricting cases routinely present computational
analyses to courts.
The opposing party can challenge such analyses on the grounds that the
expert's results are not reproducible: that running the same algorithm
on the same data would not produce the same results, or that the expert
has not disclosed the inputs and parameters used.

The \bisect{} audit chain addresses this challenge directly.
The manifest contains everything needed to reproduce the analysis:
algorithm version hash, parameter configuration, data download hashes,
adjacency construction hash, and output plan hash.
Any qualified expert can run \texttt{bisect label-verify} and confirm---
or refute---the original expert's claim.

This converts redistricting expert testimony from a credibility contest
(whose algorithm is more trustworthy?) into a verifiable computation
(does the manifest check out?).
